% CODEX: wstaw w 30b-geometrizacja-thurstona.tex bezpośrednio przed
% \section{Twierdzenie geometryzacyjne}; docelowo przed nową sekcją
% "Przykłady hiperboliczne...", jeśli oba dopiski są używane.
\begin{przyklad}[Liczba Eulera wiązki jednostkowych wektorów]
\label{prz:thur-dop-seifert-euler}
Niech $\Sigma_g$ będzie zamkniętą, zorientowaną powierzchnią
rodzaju $g\ge2$ z metryką hiperboliczną. Rzut
$\pi:U\Sigma_g\to\Sigma_g$ jest główną wiązką
$\mathrm{SO}(2)$, ponieważ w każdym punkcie wektory
jednostkowe tworzą okrąg. Jej klasa Eulera jest klasą
Eulera zorientowanej wiązki stycznej $T\Sigma_g$.

W lokalnej zorientowanej ramie ortonormalnej $e_1,e_2$
zapiszmy $\nabla e_1=A\,e_2$ oraz
$\nabla e_2=-A\,e_1$. Z równania strukturalnego
koneksji na powierzchni otrzymujemy
$\dd A=-K\,\mathrm{vol}_{\Sigma_g}$
przy tej konwencji formy $A$. Po uwzględnieniu orientacji
włókna reprezentantem klasy Eulera jest
$-\dd A/(2\pi)$, zatem
\[
 \langle e(T\Sigma_g),[\Sigma_g]\rangle
 =\frac1{2\pi}\int_{\Sigma_g}K\,\mathrm{vol}_{\Sigma_g}
 =\chi(\Sigma_g)=2-2g.
\]
Ostatnia równość jest twierdzeniem Gaussa--Bonneta
z Sekcji~\ref{sec:gauss-bonnet}, a nie osobnym postulatem
geometryzacji. W szczególności liczba Eulera
tej wiązki jest niezerowa, więc $U\Sigma_g$ nie jest
wiązką iloczynową $\Sigma_g\times S^1$.
To jest konkretny powód rozróżnienia w tabeli
\ref{thm:thurston-seifert-table} geometrii
$\widetilde{\mathrm{SL}}(2,\R)$ i $\mathbb H^2\times\R$.
\end{przyklad}
